\begin{abstract}
Diabetic Retinopathy (DR) is a leading cause of preventable blindness. Federated
learning (FL) allows hospitals to train a shared DR grading model without exchanging
retinal images, but when hospitals see very different patients, standard FedAvg can
collapse. \textbf{OrdinalFed} is studied as a remedy with two components:
(1)~an \emph{OrdinalLoss} that penalises grading errors by their squared distance,
and (2)~a \emph{QWK-weighted aggregation} in which the server scores each client
model by its Quadratic Weighted Kappa (QWK) on a 50-image server set and moves the
aggregation weights from image counts towards these scores over the rounds. On the
DDR benchmark, every configuration was repeated with three seeds and scored on the
held-out test split. Under extreme label skew (Dirichlet $\alpha=0.1$), FedAvg
collapsed in every seed (worst-round test QWK 0.378), and OrdinalFed raised the
worst round to 0.514 without any change at the hospitals and with about 1.2~s of
extra server computation per round. Ablations showed that this gain comes from
giving small hospitals more weight: equal client weights did better (0.630), while
the QWK scores and the OrdinalLoss added no measurable benefit. FedProx was the most
stable method (0.641), and adding its proximal term to OrdinalFed came close to it
(0.627). On the IID and realistic multi-site partitions, all methods were within
0.01 QWK of each other.
\end{abstract}

\begin{IEEEkeywords}
Diabetic Retinopathy, Federated Learning, Non-IID,
Ordinal Loss, QWK-Weighted Aggregation, EfficientNet-B3,
Medical Imaging
\end{IEEEkeywords}
